\section{Experiments}
\label{experiments}

\subsection{Research Questions}
\label{sec:exp-rqs}

The evaluation is designed around the following research questions:
\begin{enumerate}[leftmargin=*]
  \item RQ1: How completely does the workflow recover requirement elements from natural-language requirements?
  \item RQ2: How consistently does the workflow preserve requirement semantics in the structured representation?
  \item RQ3: How does knowledge guidance affect graph quality and representation accuracy?
  \item RQ4: How does human-in-the-loop refinement change the final requirement graph?
\end{enumerate}

\subsection{Experimental Setup}
\label{sec:exp-setup}

This subsection should report the implementation setting, agent prompts, knowledge sources, annotation process, and graph serialization format used in the study.
The setup should be sufficient for reproducing the requirement structure modeling workflow.

\subsection{Datasets}
\label{sec:exp-datasets}

The datasets should consist of natural-language requirement documents or requirement statements with sufficient domain context for semantic decomposition.
For each dataset, the final paper should report the source, domain, number of requirements, annotation protocol, and any exclusion criteria.

\subsection{Baselines}
\label{sec:exp-baselines}

The evaluation should compare the proposed workflow against focused baselines, such as direct LLM extraction, single-agent parsing, and template-only extraction.
Baselines should produce the same target representation so that comparison is made at the structure-modeling level.

\subsection{Evaluation Metrics}
\label{sec:exp-metrics}

\subsubsection{Structure Completeness}
\label{sec:exp-completeness}

Structure completeness measures whether required element types and relations are present in the generated representation.

\subsubsection{Semantic Consistency}
\label{sec:exp-consistency}

Semantic consistency measures whether graph elements preserve the meaning of the original requirement text.

\subsubsection{Graph Quality}
\label{sec:exp-graph-quality}

Graph quality evaluates relation validity, node typing, evidence traceability, and structural coherence.

\subsubsection{Requirement Representation Accuracy}
\label{sec:exp-accuracy}

Representation accuracy compares generated elements and relations against a human-reviewed reference representation.

\subsection{Main Results}
\label{sec:exp-main-results}

This subsection should report empirical results for the research questions without introducing unsupported claims.
Tables should separate element-level and relation-level performance.

\subsection{Ablation Study}
\label{sec:exp-ablation}

The ablation study should examine the contribution of knowledge guidance, role-separated agents, and human refinement.
Each ablation should preserve the same input requirements and evaluation metrics.

\subsection{Case Study}
\label{sec:exp-case-study}

The case study should trace a small set of requirements from raw text to semantic decomposition and final requirement graph.
It should highlight representation decisions, not downstream engineering tasks.
